\documentclass[10pt,a4paper]{article}
\usepackage[T1]{fontenc}
\usepackage[utf8]{inputenc}
\usepackage[margin=2cm]{geometry}
\usepackage{booktabs,longtable,array}
\usepackage{listings,xcolor}
\usepackage{enumitem}
\usepackage{hyperref}
\usepackage{amssymb}

\lstset{basicstyle=\ttfamily\small,columns=fullflexible,breaklines=true,
        frame=single,framesep=3pt,backgroundcolor=\color{black!4},
        keepspaces=true,showstringspaces=false}
\newcommand{\cmd}[1]{\texttt{\detokenize{#1}}}
\newcolumntype{P}[1]{>{\raggedright\arraybackslash}p{#1}}
\setlist{nosep}
\setlength{\parindent}{0pt}
\setlength{\parskip}{4pt}

\title{LanCall --- Phase 10 Test Plan (Essentials)\\\large Phases 4--9 on the final build}
\date{20 November 2026}
\author{}

\begin{document}
\maketitle

\textbf{Goal:} enough testing to claim every feature and fill the report's key metrics.
\begin{itemize}
  \item \textbf{Main plan:} 3 Linux laptops, \textbf{$\approx$1\,h\,45\,min}, plus a 10\,min
    wrap-up afterwards.
  \item \textbf{Windows:} in its own section at the end
    (Section~\ref{sec:windows}), run once Phase~8 can run lancall.
\end{itemize}

\section{Hardware, roles, conventions}

\begin{longtable}{@{}lllP{10.5cm}@{}}
\toprule
Id & OS & Floor & Fixed jobs\\
\midrule
L1 & Linux & A & loopback experiments, \cmd{lancall_directory} all session, group host\\
L2 & Linux & A & floor-A partner\\
L3 & Linux & B & floor-B partner, relay host, \cmd{lancall_bridge}\\
\bottomrule
\end{longtable}

\begin{itemize}
  \item There is one wall port per room.
  \item Log names: \cmd{--log logs/<test>_<laptop>.csv}.
  \item Every lancall gets \cmd{--directory <L1 ip>}.
  \item Headphones in every call. netem off after every use.
\end{itemize}

\section{Build status (main at \cmd{5b6a718}, checked 20 Nov)}

\begin{longtable}{@{}P{1.8cm}P{7.2cm}P{7.0cm}@{}}
\toprule
Phase & On main & Still missing\\
\midrule
8 Windows & platform layer and socket class on Winsock; CI on Ubuntu and Windows (libraries and their tests only) & signaling, discovery, camera, audio, and the \cmd{lancall} / \cmd{lancall_relay} / \cmd{lancall_directory} programs on Windows. Every Windows test is blocked until these land (Section~\ref{sec:windows})\\
9 UI & app core shared by \cmd{lancall} and \cmd{lancall_gui}. GUI: start screen; Peers panel (Call / Add / Watch); Call panel (incoming popup, members with safety codes, mesh / relay / multicast, Hang up); video grid; messages panel with a command box & no separate chat, stats or settings panels; chat goes through the command box. Needs \cmd{libglfw3-dev}. Logs the same rows as lancall\\
Bugs & --- & \#11 (frozen tile on an immediate rejoin) and the \cmd{av_sync_ms} outlier: not fixed, not in the NOTES.md issue table\\
\bottomrule
\end{longtable}

The terminal \cmd{lancall}'s flags, commands, \cmd{list} output and log events are unchanged,
so all the scripts work as before.

\section{Already done --- not repeated}

The tests below ran on the pre-Phase-5 build and are the Phase~4 results:
\begin{itemize}
  \item 1a self-call: audio p50 84.0\,ms;
  \item 1b group signaling;
  \item 2a 1:1 call: audio p50 88.2\,ms, 15\,fps;
  \item 2b 5\% loss: \#7 fixed, no stall;
  \item 2c bandwidth cap;
  \item 2d relay pair;
  \item 4a 3-person cross-floor mesh: audio p50 80.9\,ms, 0\% loss;
  \item 4b rejoin: works after 10\,s, \textbf{frozen tile on an immediate rejoin (\#11, open)};
  \item E6: the cross-floor penalty is about 0.
\end{itemize}

4c wasn't done; G2 covers it.

\section{Before the session (not counted)}

\begin{longtable}{@{}lP{15cm}@{}}
\toprule
\# & Task\\
\midrule
P1 & Fix or document \#11 and the \cmd{av_sync_ms} outlier from the partial 4c log. Add both to the NOTES.md issue table.\\
P2 & Done: \cmd{--no-nack} in E2, plus the helper scripts \cmd{call_timings.sh}, \cmd{exp_e7_soak.sh} and \cmd{exp_d1_discovery.sh}.\\
P3 & Aric, offline: \cmd{log_summary --per-second logs/2b_Aric.csv} ($\approx$5\% during the both-sides part closes \#9). Check that the last \cmd{video_level} in \cmd{logs/2c_Aric.csv} is 0.\\
P4 & Tag \cmd{test-final}. \textbf{Every laptop:} \cmd{sudo apt install libglfw3-dev}, build, run ctest (save the output) and devcheck, set \cmd{sudo sysctl -w net.core.rmem_max=8388608}, and check that \cmd{build/gui/lancall_gui} exists (without GLFW it is skipped silently).\\
P5 & Start \cmd{lancall_gui} from a normal terminal (Ctrl+Alt+T), not VS Code's: the snap build of VS Code makes the GUI abort (NOTES.md).\\
\bottomrule
\end{longtable}

\section{Timeline}

$T$ is minutes from the start.

\begin{longtable}{@{}lP{1.4cm}P{4.4cm}P{3.8cm}P{3.8cm}@{}}
\toprule
Slot & $T$ & L1 (A) & L2 (A) & L3 (B)\\
\midrule
S0 & 0--10 & start directory, post IPs & post IPs & post IPs\\
S1 & 10--35 & loopback: E8, E4, E3, E9, E1, D1 & 1:1 with L3: C1, C2, C5, C7, C8, C6 & 1:1 with L2\\
S2 & 35--55 & 15\,min soak with L2 & 15\,min soak & E2 (spare, loopback)\\
S3 & 55--75 & group host: E5, G4, G7, G2 & member & member, then relay host\\
S4 & 75--95 & multicast + broadcast host, E10 presenter & viewer / idle & bridge, viewer\\
S5 & 95--105 & \textbf{UI} (\cmd{lancall_gui}) & \textbf{UI} & terminal member\\
after & +10 & wrap-up & wrap-up & wrap-up\\
\bottomrule
\end{longtable}

\textbf{If behind}, cut in this order: E3 $\to$ the soak (to 10\,min) $\to$ E10 $\to$ G7.

\section{S0 --- start (10 min)}

\begin{itemize}
  \item L1: \cmd{./build/control/lancall_directory} in its own terminal, left running.
  \item Everyone: post their IP to the group chat; if ufw is on, open the ports.
\end{itemize}

\begin{lstlisting}
ip -4 addr | grep 10.8.
sudo ufw allow 5004:5016/udp && sudo ufw allow 5011/tcp     # only if ufw is on
\end{lstlisting}

\section{S1 (25 min)}

\subsection{L1 --- loopback experiments, strictly one at a time}

\begin{longtable}{@{}lP{7.0cm}lP{6.0cm}@{}}
\toprule
Step & Command & Time & Report numbers / pass\\
\midrule
E8 & \cmd{LOSSES="5 10 20" DURATION=10 tools/exp_e8_repair.sh} & 3\,min & complete frames and residual loss for none / NACK / FEC / both; upload cost\\
E4 & \cmd{tools/exp_e4_cap.sh} & 3\,min & steps down $<$5\,s after the cap, climbs back to L0 after it (\#6)\\
E3 & \cmd{JITTERS="0 20 40" BUFFERS="40 60 100" DURATION=15 tools/exp_e3_jitter.sh} & 3\,min & smallest buffer with late drops $\approx$0 for each jitter\\
E9 & \cmd{./build/tools/crypto_bench} (tee to \cmd{logs/E9_L1.txt}) & 1\,min & $\mu$s per packet to seal/open; size overhead\\
E1 & \cmd{timeout -s INT 90 ./build/call/lancall --name T --peer 127.0.0.1 --log logs/E1_L1.csv} & 1.5\,min & final latency and A/V sync (\cmd{log_summary})\\
D1 & \cmd{tools/exp_d1_discovery.sh} & 1\,min & discovery time, target $<$5\,s\\
\bottomrule
\end{longtable}

\subsection{L2 $\leftrightarrow$ L3 --- cross-floor 1:1 on the final build}

Both run \cmd{./build/call/lancall --name <me> --directory <L1 ip> --log logs/<test>_<me>.csv}.
L2 calls L3.

\begin{longtable}{@{}P{2.2cm}P{6.4cm}lP{6.1cm}@{}}
\toprule
Test & Steps & Time & Report numbers / pass\\
\midrule
C1 clean & call, 2\,min, \cmd{stats}, \cmd{members} & 3\,min & ``encrypted''; \textbf{same safety code on both}; audio/video p50, fps, rtt (\cmd{log_summary}); setup time via \cmd{tools/call_timings.sh}: accept $\to$ first frame $<$2\,s\\
C2 5\% loss & \cmd{tools/netem.sh on 5} on both, 2\,min, \cmd{stats} at 1\,min, \cmd{off} & 3\,min & measured loss $\approx$5\%, residual after NACK $\approx$0, \cmd{nack N}$>$0, no stall\\
C5 outage & clean call; L3 runs the listing below twice, 20\,s apart & 2\,min & \cmd{call_timings.sh logs/C5_L2.csv}: recovery (gap $-$ 1\,s) $<$2\,s\\
C7 chat & \cmd{text L3 hi} outside a call; \cmd{text all ...} in a call; to an offline name; a 501-byte message; emoji & 3\,min & delivered, or a clear error\\
C8 old client refused & \cmd{nc <L2 ip> 5011}, then type \cmd{INVITE} & 1\,min & answered \cmd{DECLINE reason=old-version}\\
C6 Ctrl+C & Ctrl+C on L3 mid-call & 1\,min & L2's call ends cleanly; calling again works\\
\bottomrule
\end{longtable}

\begin{lstlisting}
tools/netem.sh on 100; sleep 1; tools/netem.sh off      # C5: one 1 s outage
\end{lstlisting}

\section{S2 (20 min)}

\subsection{L1 $\leftrightarrow$ L2 --- E7 soak (15 min)}

\begin{lstlisting}
WARMUP_S=60 WINDOW_S=180 MIN_CALL_S=840 tools/exp_e7_soak.sh run L1    # L2 likewise
\end{lstlisting}

\begin{itemize}
  \item L1 calls L2. Don't touch it for 15\,min; then \cmd{hangup}, then \cmd{quit}.
  \item \textbf{Pass:} from the printed summary:
    \begin{itemize}
      \item memory growth $<$10\%;
      \item audio p50 moved $<$5\,ms;
      \item fps steady.
    \end{itemize}
  \item The mean CPU is the encrypted call's CPU (E9).
\end{itemize}

\subsection{L3 --- E2 (spare)}

\cmd{DURATION=10 tools/exp_e2_loss.sh}: whole frames vs 10 strips at 0--20\% loss, on the
final build, with NACK off.

\section{S3 (20 min) --- group calls}

\subsection*{E5 --- mesh vs relay (2 runs, $\sim$5 min)}

\begin{lstlisting}
# members (L2, L3)
EXTRA="--directory <L1 ip>" tools/exp_e5_group.sh run <mode> 3 <name>
# host L1
EXTRA="--directory <L1 ip>" RELAY=me tools/exp_e5_group.sh run <mode> 3 L1 L2 L3
\end{lstlisting}

\begin{itemize}
  \item Runs: \textbf{mesh 3, relay 3}. The 2-person numbers come from C1.
  \item \textbf{Report:} upload per sender (mesh $\approx$ 2 copies, relay 1), latency in each
    mode, 15\,fps. Collect into \cmd{results/e5_group.csv}.
\end{itemize}

\subsection*{One call --- L1 hosts L2 and L3 ($\sim$15 min)}

\begin{longtable}{@{}P{2.8cm}P{7.0cm}P{6.4cm}@{}}
\toprule
Test & Steps & Pass\\
\midrule
G4 group chat + codes & L2 \cmd{text all hi}; L2 \cmd{text L3 psst}; \cmd{members} on all & delivered as addressed; safety codes match per pair\\
G7 dead member & pull L3's cable for 20\,s, plug it back, L1 \cmd{add L3} & the others drop L3 after $\sim$10\,s; the call carries on; the re-add works\\
G2a relay in the call & L1 \cmd{relay me}, 1\,min, \cmd{mesh} & media continues; \cmd{call_timings.sh} gaps $<$2\,s; \cmd{members} agrees\\
G2b relay on a non-member (replaces 4c) & L3 hangs up and runs \cmd{./build/net/lancall_relay}; L1 \cmd{relay <L3 ip>} with L2, 1\,min, then \cmd{mesh} & media continues through L3; the relay prints 2 members; \cmd{av_sync_ms} min is sane (\cmd{log_summary})\\
G6 rejoin (\#11) & \textbf{only if P1 fixed it:} L3 \cmd{hangup}, L1 \cmd{add L3} at once & L3's tile is live on L2\\
\bottomrule
\end{longtable}

\section{S4 (20 min) --- multicast, broadcast, E10}

\begin{longtable}{@{}P{2.8cm}P{7.0cm}P{6.4cm}@{}}
\toprule
Test & Steps & Pass\\
\midrule
M1 multicast call + auto switch & L1 calls L2, \cmd{multicast}, 1\,min, \cmd{stats}; then \cmd{add L3} & \cmd{multicast 239.255.77.x} with one upload copy; after the add, \textbf{switches to a relay on its own}; L3 gets media (\cmd{call_timings.sh} for the gap)\\
B1 broadcast & L1 \cmd{broadcast demo}; L2 \cmd{watch L1} 1\,min, \cmd{unwatch}; L1 tries \cmd{call L3} & L2 sees and hears it; the call is refused while broadcasting\\
B3 bridge & L1 (still broadcasting) \cmd{bridge <L3 ip>}; L3 runs the printed \cmd{lancall_bridge} line, then \cmd{watch L1-bridge} in its own lancall & video plays on floor B via the bridge\\
\bottomrule
\end{longtable}

\subsection*{E10 --- reduced (3 runs $\times$ $\sim$2 min)}

The commands are at the top of \cmd{tools/exp_e10_multicast.sh}.

\begin{longtable}{@{}llll@{}}
\toprule
Run & L1 & L2 & L3\\
\midrule
unicast 2 & presenter & viewer & viewer\\
multicast 1 & presenter & viewer & --\\
multicast idle & presenter, nobody watching & idle & idle (other floor)\\
\bottomrule
\end{longtable}

\textbf{Report:}
\begin{itemize}
  \item presenter upload: unicast with 2 viewers vs multicast (the same for any number of
    viewers);
  \item the flood cost on a same-floor laptop that isn't watching (L2) against the other
    floor (L3).
\end{itemize}

\section{S5 (10 min) --- Phase 9 UI gate, on Linux}

L1 and L2 run \cmd{lancall_gui --name <me> --directory <L1 ip> --log logs/G9_<me>.csv}; L3
uses the terminal.
\begin{itemize}
  \item \textbf{G9.1:} using only the UI on L1:
    \begin{itemize}
      \item call L2 from the Peers panel (L2 accepts from the popup), then add L3;
      \item chat by typing \cmd{text all ...} in the command box (there's no chat panel);
      \item switch to ``Relay on this laptop'' and back to Mesh;
      \item check the safety codes in the members table, then hang up.
    \end{itemize}
  \item \textbf{G9.2:} L1 types \cmd{broadcast ui} in the command box; L2 watches with the
    Peers panel's Watch button.
  \item \textbf{G9.4:} a 2-minute L1$\leftrightarrow$L2 call in the UI. Compare its
    audio/video p50 (\cmd{log_summary}) with the soak's first minutes, and note CPU from
    \cmd{top}. Pass: within $\sim$5\,ms.
\end{itemize}

\section{Wrap-up (10 min, after testing)}

\begin{itemize}
  \item \cmd{tools/netem.sh off} on every Linux laptop.
  \item Copy the experiment CSVs and key logs to \cmd{results/}; add headline rows to
    \cmd{summary.csv}.
\end{itemize}

\section{Report metrics and features}

\begin{longtable}{@{}P{5.4cm}P{4.4cm}P{6.0cm}@{}}
\toprule
Metric & Target & Test\\
\midrule
Latency & $<$300\,ms (aim $<$150) & E1, C1, E5\\
Loss tolerance and repair & usable at 5\% & C2, E8, E2\\
A/V sync & $<$150\,ms & E1, C1\\
Frame rate & $\ge$15\,fps at 480p & C1, E5\\
Group size & $\ge$4 & E5 shows 3 here; 4 needs the Windows laptop (W5)\\
Call setup & $<$2\,s & C1\\
Discovery & $<$5\,s & D1\\
Recovery after a 1\,s outage & $<$2\,s & C5\\
Adaptive under a cap & down $<$5\,s, back up & E4\\
Jitter vs buffer & sweet spot & E3\\
Encryption cost & $\mu$s, \%, CPU & E9, E7\\
Upload: mesh / relay / multicast & per sender & E5, E10\\
Stability & 15\,min soak & E7\\
UI overhead & $\le$5\,ms & G9.4\\
\bottomrule
\end{longtable}

\textbf{Features claimed $\to$ test:}
\begin{itemize}[label=$\square$]
  \item Discovery: D1. Directory: the S1 cross-floor list. Old clients refused: C8.
  \item Calls: C1, C6. Group mesh, relay, \cmd{relay me} and a relay on a non-member: E5, G2.
    Dead member: G7.
  \item NACK: C2. FEC: E8. Adaptive video: E4. Chat (1:1 and group): C7, G4.
  \item Encryption and safety codes: C1, G4.
  \item Broadcast and watch: B1. Bridge: B3. Multicast call and auto switch: M1.
  \item UI: G9 (Linux).
\end{itemize}

\textbf{Cut from the full plan} (not claimed as separately tested):
\begin{itemize}
  \item C3 one-sided loss (\#9 is settled offline in P3).
  \item C4 FEC inside a call (E8 measures FEC).
  \item Most of C6's life-cycle cases, C9 directory restart, and the old-binary build (C8 uses
    \cmd{nc}).
  \item E5 with 2 people.
  \item Busy callee, host hangup, multicast refused, two broadcasters at once.
  \item E10 relay runs.
  \item UI settings.
  \item 30\,min soak, now 15\,min.
\end{itemize}

\section{Windows (Phase 8) --- run once \texttt{lancall.exe} works}\label{sec:windows}

\textbf{Prerequisite:} on W, build the tag with MSYS2 UCRT64 from the documented steps; ctest
passes apart from the Linux-only tests (gate W1). Add a firewall rule (or mark the network
Private). Install a GUI build if Phase 9 builds on Windows.

\textbf{Where:} W in a floor-B room with L3. L1 and L2 are needed only for W5 and W6.

\textbf{Time:} about 25\,min as a separate session.
\begin{itemize}
  \item \textbf{Alternative, if W is ready on the day:} fold it into the main plan, which
    adds only about 10\,min ($\approx$1\,h\,55\,min in total).
  \item \textbf{How to fold it in:}
    \begin{itemize}
      \item W2--W4 run on L3$\leftrightarrow$W during S2, alongside the soak; L3 skips E2.
      \item W5 adds mesh 4 and relay 4 to E5, and runs G2b with \cmd{lancall_relay.exe} on W,
        instead of on L3.
      \item W6 joins S5.
    \end{itemize}
\end{itemize}

\begin{longtable}{@{}P{2.6cm}P{6.8cm}P{6.8cm}@{}}
\toprule
Gate & Steps & Pass\\
\midrule
W1 build & clean clone, documented steps, ctest & builds; tests pass apart from the Linux-only ones\\
W2 discovery + 1:1 & \cmd{list} on W and L3; W calls L3 for 2\,min, then L3 calls W; chat both ways; \cmd{members} & each lists the other within 5\,s (the subnet broadcast address on Windows); audio and video both ways; encrypted; same safety codes; no audio crackle (\cmd{timeBeginPeriod}); L3's log gives audio p50\\
W3 watch & L3 \cmd{broadcast floorB}; W \cmd{watch L3} 1\,min, \cmd{unwatch} & video and sound on W; \cmd{list} shows the title\\
W4 multicast call & L3 calls W, then \cmd{multicast} & the status shows \cmd{multicast 239.255.77.x} on both; media flows\\
W5 group + relay on Windows & E5 \textbf{mesh 4, relay 4} (W accepts by hand). Then L1 hosts L2 and L3, and W runs \cmd{lancall_relay.exe}; L1 \cmd{relay <W ip>}, 1\,min, \cmd{mesh} & W's tile is live and W sees 3 tiles; the 4-person runs give the ``$\ge$4 people'' metric; the relay on Windows forwards for 3 members\\
W6 UI on Windows & W runs \cmd{lancall_gui}; L1 (GUI) adds W; W accepts from the popup; W watches L3's broadcast with the Watch button & the session works on Windows, from the UI only\\
W7 directory on Windows & W runs \cmd{lancall_directory.exe}; L3 restarts with \cmd{--directory <W ip>} & peers from floor A show up on L3\\
\bottomrule
\end{longtable}

\textbf{Report rows:}
\begin{itemize}
  \item W2's latency (Windows $\leftrightarrow$ Linux);
  \item W5's 4-person E5 rows;
  \item pass/fail for W1--W7.
\end{itemize}
If Phase~8 never lands, the report states that Windows support covers the platform layer and
CI only.

\end{document}
